\section*{Ex. 2.5 (20 pts) --- Pseudo-code for \texttt{fork}, \texttt{exec}, \texttt{waitpid}, \texttt{exit}}

Following the hint, each algorithm below is the \emph{kernel-side implementation} of the
call: the code the operating system runs when a user process traps into it. All four
operate on the process table, whose entries are process control blocks (PCBs);
\textsc{Current} denotes the PCB of the calling process, and a negative return value is
the usual \texttt{errno} convention.

\begin{algorithm}[H]
\caption{\texttt{fork()} --- create a copy of the calling process}
\label{p4:alg:fork}
\begin{algorithmic}[1]
\Procedure{Fork}{\,}
  \If{no free process-table slot \textbf{or} no free pid}
    \State \Return $-$\texttt{EAGAIN}
  \EndIf
  \State $c \gets$ \Call{AllocatePCB}{\,};\quad $c.\mathit{pid} \gets$ \Call{AllocatePid}{\,}
  \State copy from \textsc{Current} into $c$: uid, gid, cwd, root, umask,
         \Statex \hspace{\algorithmicindent} signal dispositions, resource limits, priority
  \State $c.\mathit{parent} \gets \textsc{Current}$; append $c$ to $\textsc{Current}.\mathit{children}$
  \State $c.\mathit{addrspace} \gets$ \Call{CopyOnWriteShare}{$\textsc{Current}.\mathit{addrspace}$}
  \State mark every writable page read-only in both processes
         \Comment{copy on first write}
  \ForAll{descriptors $fd$ in $\textsc{Current}.\mathit{fdtable}$}
    \State $c.\mathit{fdtable}[fd] \gets \textsc{Current}.\mathit{fdtable}[fd]$
    \State increment the reference count of that open-file entry
  \EndFor
  \State $c.\mathit{cputime} \gets 0$;\quad $c.\mathit{pending\ signals} \gets \emptyset$
  \State $c.\mathit{savedregs} \gets \textsc{Current}.\mathit{savedregs}$
         \Comment{child resumes at the instruction after the trap}
  \State $c.\mathit{savedregs}.\mathit{retval} \gets 0$
         \Comment{what the child will see \texttt{fork} return}
  \State $c.\mathit{state} \gets \textsc{Ready}$; \Call{Enqueue}{\emph{ready queue}, $c$}
  \State \Return $c.\mathit{pid}$
         \Comment{what the parent will see \texttt{fork} return}
\EndProcedure
\end{algorithmic}
\end{algorithm}

The two return-value assignments --- $0$ into the child's saved register, the child's pid
into the parent's --- are the whole ``\texttt{fork} returns twice'' phenomenon: one call
site, two saved register sets, two different values in the return register.

\begin{algorithm}[H]
\caption{\texttt{exec(path, argv, envp)} --- replace the caller's program image}
\label{p4:alg:exec}
\begin{algorithmic}[1]
\Procedure{Exec}{$\mathit{path}, \mathit{argv}, \mathit{envp}$}
  \State $\mathit{inode} \gets$ \Call{Namei}{$\mathit{path}$};
         \textbf{if} $\mathit{inode} = \textsc{Null}$ \textbf{then return} $-$\texttt{ENOENT}
  \State \textbf{if} $\mathit{inode}$ is not executable by $\textsc{Current}.\mathit{uid}$
         \textbf{then return} $-$\texttt{EACCES}
  \State $\mathit{hdr} \gets$ \Call{ReadHeader}{$\mathit{inode}$}
  \State $\mathit{loader} \gets$ \Call{LoaderFor}{$\mathit{hdr}.\mathit{magic}$}
         \Comment{magic number selects a.out / ELF / script}
  \State \textbf{if} $\mathit{loader} = \textsc{Null}$ \textbf{then return} $-$\texttt{ENOEXEC}
  \State $\mathit{kargv}, \mathit{kenvp} \gets$ \Call{CopyIntoKernel}{$\mathit{argv}, \mathit{envp}$}
         \Comment{they live in the old address space}
  \Statex \hspace{\algorithmicindent}\emph{--- point of no return: every error above is still
          reportable to the caller ---}
  \State \Call{DestroyAddressSpace}{$\textsc{Current}.\mathit{addrspace}$}
  \State $\textsc{Current}.\mathit{addrspace} \gets \mathit{loader}.$\Call{MapImage}{$\mathit{inode}$}
         \Comment{text r-x, data rw, bss zeroed}
  \State create a fresh heap and user stack
  \ForAll{signals $s$ caught by a user handler}
    \State disposition of $s \gets \textsc{Default}$
           \Comment{the handler's code is gone}
  \EndFor
  \Statex \hspace{\algorithmicindent}\emph{ignored dispositions and the blocked signal
          mask are left untouched}
  \State \textbf{close} every descriptor marked close-on-exec
  \State \textbf{if} $\mathit{inode}$ is set-uid \textbf{then}
         $\textsc{Current}.\mathit{euid} \gets \mathit{inode}.\mathit{owner}$
  \State $\mathit{sp} \gets$ \Call{PushOntoNewStack}{$\mathit{kargv}, \mathit{kenvp}$}
  \State $\textsc{Current}.\mathit{savedregs} \gets \{\,\mathit{pc}: \mathit{hdr}.\mathit{entry},\;
         \mathit{sp}: \mathit{sp},\; \text{others}: 0\,\}$
  \State \textbf{return to user mode}
         \Comment{on success there is no caller left to return \emph{to}}
\EndProcedure
\end{algorithmic}
\end{algorithm}

\begin{algorithm}[H]
\caption{\texttt{waitpid(pid, statusp, options)} --- collect a terminated child}
\label{p4:alg:waitpid}
\begin{algorithmic}[1]
\Procedure{Waitpid}{$\mathit{pid}, \mathit{statusp}, \mathit{options}$}
  \Loop
    \State $\mathit{found} \gets \textsc{False}$
    \ForAll{$c \in \textsc{Current}.\mathit{children}$ matching $\mathit{pid}$}
      \Comment{$\mathit{pid}=-1$ matches any child}
      \State $\mathit{found} \gets \textsc{True}$
      \If{$c.\mathit{state} = \textsc{Zombie}$}
        \State \textbf{if} $\mathit{statusp} \neq \textsc{Null}$ \textbf{then}
               \Call{CopyOut}{$\mathit{statusp}, c.\mathit{exitstatus}$}
        \State $\textsc{Current}.\mathit{childcputime} \mathrel{+}= c.\mathit{cputime} + c.\mathit{childcputime}$
        \State remove $c$ from $\textsc{Current}.\mathit{children}$
        \State $p \gets c.\mathit{pid}$;\quad \Call{FreePCB}{$c$};\quad \Call{FreePid}{$p$}
               \Comment{the reaping step}
        \State \Return $p$
      \EndIf
    \EndFor
    \State \textbf{if} $\mathit{found} = \textsc{False}$ \textbf{then return} $-$\texttt{ECHILD}
    \State \textbf{if} $\mathit{options}$ contains \texttt{WNOHANG} \textbf{then return} $0$
    \State \textbf{if} a signal is pending for \textsc{Current} \textbf{then return} $-$\texttt{EINTR}
    \State \Call{Sleep}{$\textsc{Current}$, \emph{its wait channel}}
           \Comment{a child's \texttt{exit} wakes us; then rescan}
  \EndLoop
\EndProcedure
\end{algorithmic}
\end{algorithm}

\begin{algorithm}[H]
\caption{\texttt{exit(status)} --- terminate the calling process}
\label{p4:alg:exit}
\begin{algorithmic}[1]
\Procedure{Exit}{$\mathit{status}$}
  \ForAll{open descriptors $fd$ of \textsc{Current}}
    \State \Call{Close}{$fd$}
           \Comment{drops the reference count; the last close releases the file}
  \EndFor
  \State \Call{ReleaseAddressSpace}{$\textsc{Current}.\mathit{addrspace}$}
  \State release cwd, root, timers and file locks
  \ForAll{$c \in \textsc{Current}.\mathit{children}$}
    \State $c.\mathit{parent} \gets \mathit{init}$; move $c$ to $\mathit{init}.\mathit{children}$
    \State \textbf{if} $c.\mathit{state} = \textsc{Zombie}$ \textbf{then}
           \Call{Signal}{$\mathit{init}$, \texttt{SIGCHLD}}; \Call{Wakeup}{$\mathit{init}$}
           \Comment{so \emph{init} reaps it}
  \EndFor
  \State $\textsc{Current}.\mathit{exitstatus} \gets \mathit{status}$; record CPU time and accounting in the PCB
  \State $\textsc{Current}.\mathit{state} \gets \textsc{Zombie}$
         \Comment{PCB and pid stay allocated until reaped}
  \State \Call{Signal}{$\textsc{Current}.\mathit{parent}$, \texttt{SIGCHLD}};
         \Call{Wakeup}{$\textsc{Current}.\mathit{parent}$}
  \State \Call{Schedule}{\,}
         \Comment{picks another process; \texttt{exit} never returns}
\EndProcedure
\end{algorithmic}
\end{algorithm}

\subsection*{How the four calls are inter-related}

They are not four independent services but two pairs that bracket the life of a process.

\texttt{fork} is the only way a Unix system makes a new process, and \texttt{exec} is the
only way a process changes the program it runs. Neither alone can start a new program:
\texttt{fork} produces a second copy of the caller, \texttt{exec} replaces the caller's
own image. Running a new program is therefore necessarily both of them, \texttt{fork}
then \texttt{exec} in the child. The two algorithms are built to meet: \texttt{fork}
gives the child a duplicate descriptor table and the parent's credentials, and
\texttt{exec} keeps exactly that table, less the close-on-exec entries, while discarding
everything else.

The split is what makes the arrangement useful. Between \texttt{fork} returning $0$ and
\texttt{exec} overwriting the image there is a window in which the child is still running
the parent's code and can rearrange precisely the state that \texttt{exec} preserves.
That window is where a shell implements \texttt{ls > out}: the child closes descriptor
$1$, reopens it on the file, then execs \texttt{ls}, which need know nothing about
redirection.

\texttt{exit} and \texttt{waitpid} are the matching pair at the other end, and they are a
pair for one reason: the exit status has to outlive the process that produced it.
\texttt{exit} can free the address space and the descriptors, but it cannot free the
PCB, because nobody has read the status out of it yet; so it stores the status there,
marks the process \textsc{Zombie} and signals the parent. \texttt{waitpid} is the other
half of that transaction: it finds the zombie child, copies the status out, charges its
CPU time to the parent, and only then releases the PCB and the pid. Until then the pid
stays reserved, so ``wait for process $p$'' stays unambiguous.

A shell's main loop is the whole mechanism in one cycle: \texttt{fork}; the child
\texttt{exec}s the command; the parent blocks in \texttt{waitpid}; the child
\texttt{exit}s and becomes a zombie; the parent's \texttt{waitpid} returns with the
status, frees the PCB, and prints the next prompt.

\section*{Ex. 2.6 (5 pts) --- Zombie processes}

A zombie is a process that has already run \texttt{exit()} but has not yet been reaped.
Its address space, open files and kernel stack are gone; all that is left is its
process-table entry, holding the exit status, the pid and the CPU time it used.

The kernel keeps that entry because the status belongs to the parent, and the child
cannot report it once it no longer exists. The pid must stay allocated for the same
reason: reusing it while the parent can still name it in \texttt{waitpid} would let the
parent wait on an unrelated process.

The cost is one process-table slot and one pid --- no memory, no CPU time, no
descriptors. A zombie also cannot be killed, since it is already dead: \texttt{SIGKILL}
is discarded, as no execution context is left to receive it. In \texttt{ps} it appears in
state \texttt{Z}, with its command shown as \texttt{<defunct>}.

It disappears when the parent calls \texttt{wait} or \texttt{waitpid}, or when the parent
itself exits and \emph{init} inherits and reaps it. A long-lived parent that forks
children and never waits leaks table entries until \texttt{fork} fails with
\texttt{EAGAIN}, which is why such a parent should handle or explicitly ignore
\texttt{SIGCHLD}.

An orphan is the complement: a child still \emph{alive} when its parent dies.
\emph{init} adopts it and reaps it when it eventually exits.
